% UTF-8 LaTeX fragment: reference answers and solutions

\sectionbanner{一、单选题}{参考答案与解析}

\begin{solution}{1}{C}{8.0}{代数与方程应用}
\analysis{以总量为同一基准倒推两次给出后的份额，最后检查给出张数均为整数。}
\begin{enumerate}[label=\textbf{Step \arabic*:},leftmargin=*]
  \item 设甲最初有\(a\)张，则乙最初有\(120-a\)张。
  \item 甲给出四分之一后，甲有\(\dfrac{3a}{4}\)，乙有
  \[
    120-a+\dfrac{a}{4}=120-\dfrac{3a}{4}.
  \]
  \item 乙再给出此时的三分之一后，甲最终有
  \[
    \dfrac{3a}{4}+\dfrac{1}{3}\Bigl(120-\dfrac{3a}{4}\Bigr)=\dfrac{a}{2}+40,
  \]
  乙最终有
  \[
    \dfrac{2}{3}\Bigl(120-\dfrac{3a}{4}\Bigr)=80-\dfrac{a}{2}.
  \]
  \item 令两者相等得\(a=40\)，故乙最初有\(120-40=80\)张，选C。
  \item 正向核对：最初\((40,80)\)，第一次后\((30,90)\)，第二次后\((60,60)\)；两次给出的\(10\)、\(30\)张都是整数。
\end{enumerate}
\methodsummary{先后分段基数与整张回代。}
\end{solution}

\begin{solution}{2}{B}{4.0}{排列组合与概率}
\analysis{蚂蚁从\(A\)到相对顶点\(B\)恰走\(3\)条棱，必须分别走长、宽、高三个方向各一次；用分步乘法原理计数。}
\begin{enumerate}[label=\textbf{Step \arabic*:},leftmargin=*]
  \item \(A\)与\(B\)是长方体的一对空间对角顶点，三条棱恰好覆盖三个不同方向。
  \item 第一步从\(A\)出发有\(3\)个方向可选；第二步剩\(2\)个方向；第三步只剩\(1\)个方向。
  \item 总路径数为\(3\times2\times1=6\)，选B。
\end{enumerate}
\methodsummary{三维路径先降维成方向排列，再乘积计数。}
\end{solution}

\begin{solution}{3}{C}{8.0}{排列组合与概率}
\analysis{先算全部相邻座位对，再按空空、空坐、坐坐三类分账。}
\begin{enumerate}[label=\textbf{Step \arabic*:},leftmargin=*]
  \item 全部座位的相邻对：横向\(5\times6=30\)对，纵向\(4\times7=28\)对，共\(58\)对。
  \item 两个角上空座各连\(2\)个邻座，三个内部空座各连\(4\)个，从空座端点数相邻关系共
  \[
    2\times2+3\times4=16
  \]
  次。
  \item 已知\(12\)对空坐各计一次；设空空相邻有\(k\)对，它们各被计两次，故\(12+2k=16\)，得\(k=2\)。
  \item 按三类无重无漏分配：所求坐坐相邻对数为
  \[
    58-2-12=44,
  \]
  选C。可行座次示例：空第\(1\)排\(1\)号、第\(5\)排\(7\)号以及第\(3\)排\(3,4,5\)号。
\end{enumerate}
\methodsummary{边界端点双计与三类边分账。}
\end{solution}

\begin{solution}{4}{B}{6.0}{代数与方程应用}
\analysis{原七卡是公差\(3\)的等差列；利用删中间卡后两端三卡的差先定位删位，再用总和反推首项。}
\begin{enumerate}[label=\textbf{Step \arabic*:},leftmargin=*]
  \item 若取走正中间第\(4\)张，留下的最大三张与最小三张对应项各相差\(4\times3=12\)，两组和之差恰为\(3\times12=36\)。取走第\(3\)或第\(5\)张则差变为\(33\)，取走第\(2\)或第\(6\)张则差为\(30\)。故缺项必在中间。
  \item 等差数列左右各三张关于中间编号成对，留下六张的和等于中间编号的\(6\)倍，所以中间编号为\(90\div6=15\)，首张为
  \[
    15-3\times3=6.
  \]
  \item 原七张为\(6,9,12,15,18,21,24\)。去中间后最小三张和\(27\)、最大三张和\(63\)，差\(36\)，六张总和\(90\)，选B。
\end{enumerate}
\methodsummary{等差缺项与两端差反推。}
\end{solution}

\begin{solution}{5}{B}{6.0}{几何与空间}
\analysis{把阴影面积看作\(1\)份，由长方形与甲、乙的倍数关系逐段求份数，再用乙的实际面积定出一份。}
\begin{enumerate}[label=\textbf{Step \arabic*:},leftmargin=*]
  \item 阴影为\(1\)份，则长方形\(ABCD\)为\(6\)份。甲与阴影拼成长方形，故甲为\(6-1=5\)份。
  \item 甲、乙面积和为阴影的\(9\)倍，即甲\(+\)乙\(=9\)份，因此乙为\(9-5=4\)份。
  \item 乙实际面积\(12\,\mathrm{cm}^2\)，故一份为\(12\div4=3\,\mathrm{cm}^2\)。阴影恰为一份，面积为\(3\,\mathrm{cm}^2\)，选B。
\end{enumerate}
\methodsummary{重叠倍数题优先设份数，再由拼凑关系还原各区域。}
\end{solution}

\begin{solution}{6}{B}{8.0}{代数与方程应用}
\analysis{两次试装都从原状开始，先相减找到共同袋数，再解非负整数方程。}
\begin{enumerate}[label=\textbf{Step \arabic*:},leftmargin=*]
  \item 第二次比第一次每袋多\(3\)块，总差\(143-98=45\)，故袋数为\(15\)，原有积木\(98-15=83\)块。
  \item 记小、中、大袋数为\(a,b,c\)，则
  \[
    a+b+c=15,\qquad 3a+5b+8c=83.
  \]
  消去\(a\)得\(2b+5c=38\)。
  \item \(c\)须为正偶数：\(c=2\)时\(b=14,a=-1\)不可行；\(c=4\)时\(b=9,a=2\)均正。故最小大袋数为\(4\)，选B。
\end{enumerate}
\methodsummary{权重消共同项与整数边界。}
\end{solution}

\begin{solution}{7}{D}{4.0}{代数与方程应用}
\analysis{目标和\(180\)迫使最大乘数\(127\)必须乘\(1\)，从而筛掉一半起点后再验证。}
\begin{enumerate}[label=\textbf{Step \arabic*:},leftmargin=*]
  \item 逆时针顶点取值：\(P=1,Q=1,R=0,S=1,T=0,U=1\)；乘数依次为\(3,7,15,31,63,127\)。
  \item 若最后一项乘\(0\)，前五项全乘\(1\)也只有\(119<180\)，故终点取值必为\(1\)。
  \item 从\(T\)出发逆时针顺序为\(T,U,P,Q,R,S\)，取值\(0,1,1,1,0,1\)，得
  \[
    0\cdot3+1\cdot7+1\cdot15+1\cdot31+0\cdot63+1\cdot127=180.
  \]
  故从\(T\)开始，选D。
\end{enumerate}
\methodsummary{极值估算法先锁定最大乘数对应的数字。}
\end{solution}

\begin{solution}{8}{D}{6.0}{新定义与抽象建模}
\analysis{先锁定四个互异数只能是\(9+10+12+13=44\)，再以唯一的\(12\)为锚点，按一字型与\(\mathrm{T}\)型分类清点。}
\begin{enumerate}[label=\textbf{Step \arabic*:},leftmargin=*]
  \item 表中只有\(9,10,12,13\)四种数，四数互异且和为\(44\)仅有\(9+10+12+13\)。每种合法选法恰含这四个数，且必含中心格\(12\)。
  \item 一字型（四格共线）仅横、竖两向，穿过\(12\)共\(4\)种：顶部竖条、中行左半横条、中行右半横条、正下方竖条。
  \item \(\mathrm{T}\)型为“三格一排\(+\)中间凸出一格”，在图中逐向检查得\(7\)种（含题干示例）。
  \item 合计\(4+7=11\)种，选D。
\end{enumerate}
\methodsummary{锁数—定锚—分形数：形状分类后逐类清点，不重不漏。}
\end{solution}

\begin{solution}{9}{C}{8.0}{几何与空间}
\analysis{照片只能左右平移。接触不到的盲区等于最左与最右极端位置处四个角落空白，恰为外接正方形减去相片。}
\begin{enumerate}[label=\textbf{Step \arabic*:},leftmargin=*]
  \item \(M,N\)始终贴在上下边框，竖直高度固定为相框宽\(10\,\mathrm{cm}\)。倾斜相片的外接正方向框边长亦为\(10\,\mathrm{cm}\)。
  \item 滑到最左、最右时留下的四个角落空白，合起来恰等于此外接正方形内的四个直角三角形。
  \item 盲区面积为
  \[
    10\times10-8\times8=100-64=36\,\mathrm{cm}^2,
  \]
  选C。
\end{enumerate}
\methodsummary{平移扫过面积找极端位置，用外接框面积差转化角落盲区。}
\end{solution}

\begin{solution}{10}{D}{8.0}{排列组合与概率}
\analysis{临时为六个珠位编号，最后按旋转轨道大小还原手环；镜像不合并。}
\begin{enumerate}[label=\textbf{Step \arabic*:},leftmargin=*]
  \item 红、蓝、绿各两颗的线性排布数为
  \[
    \frac{6!}{2\cdot2\cdot2}=90.
  \]
  \item 通常一个手环对应\(6\)种编号排布。例外是“三色各一颗”组成的三珠块重复两遍：此类编号排布共\(3!=6\)种，每类仅有\(3\)个不同起点，计\(6/3=2\)类。
  \item 其余\(90-6=84\)种每六个为一类，计\(14\)类。总计\(14+2=16\)类，选D。
\end{enumerate}
\methodsummary{圆周定锚与重复轨道。}
\end{solution}

\sectionbanner{二、填空题}{参考答案与解析}

\begin{solution}{11}{250}{4.0}{代数与方程应用}
\analysis{用袋数上限和铅笔总量得到费用下界，再检验整数取等。}
\begin{enumerate}[label=\textbf{Step \arabic*:},leftmargin=*]
  \item 设甲\(x\)袋、乙\(y\)袋，则
  \[
    x\ge1,\ y\ge1,\ x+y\le22,\ 12x+16y\ge288.
  \]
  化简得\(3x+4y\ge72\)。
  \item 费用\(10x+15y=5(2x+3y)\)。在约束下取\(x=16\)、\(y=6\)：袋数恰\(22\)，铅笔恰\(288\)支，费用
  \[
    10\times16+15\times6=250
  \]
  元，两种袋子均有，故下界可达且答案唯一。
\end{enumerate}
\methodsummary{资源限额与取等构造。}
\end{solution}

\begin{solution}{12}{14}{6.0}{几何与空间}
\analysis{用线段平移法把折线\(A\to C\)化为\(AB+BC\)；再利用时间相同、路程比等于速度比。}
\begin{enumerate}[label=\textbf{Step \arabic*:},leftmargin=*]
  \item 折线\(A\to C\)的水平段下移、竖直段左移后，总长\(L=AB+BC\)。
  \item 小李路程为\(L+CD=L+2\)；小王路程为\(AB+BD=L-2\)。
  \item 速度比为\(3:4\)，同时出发则路程比亦为\(3:4\)：
  \[
    \frac{L-2}{L+2}=\frac{3}{4}\implies L=14.
  \]
\end{enumerate}
\methodsummary{折线平移化直角边，再结合行程份数差量求解。}
\end{solution}

\begin{solution}{13}{55}{6.0}{几何与空间}
\analysis{把复合角两式相减消去共享小角，得到另两角的差；结合比值再用平角收尾。}
\begin{enumerate}[label=\textbf{Step \arabic*:},leftmargin=*]
  \item 记\(\angle AOB=a\)，\(\angle BOC=b\)，\(\angle COD=c\)，\(\angle DOE=d\)。则
  \[
    a+b=80^\circ,\qquad b+c=95^\circ,\qquad a:c=2:3.
  \]
  \item 两式相减得\(c-a=15^\circ\)。由\(a:c=2:3\)得\(a=30^\circ\)、\(c=45^\circ\)，进而\(b=50^\circ\)。
  \item \(OA\)与\(OE\)相反，半平角总和\(180^\circ\)，故
  \[
    d=180^\circ-(30^\circ+50^\circ+45^\circ)=55^\circ.
  \]
\end{enumerate}
\methodsummary{复合角差消共享项。}
\end{solution}

\begin{solution}{14}{32}{8.0}{几何与空间}
\analysis{先用折线的\(45^\circ\)关系识别翻折形状与原纸的交集。}
\begin{enumerate}[label=\textbf{Step \arabic*:},leftmargin=*]
  \item 折痕\(E(10,0)\)、\(F(2,8)\)与纸边成\(45^\circ\)。翻折左侧梯形面积为
  \[
    \frac{(10+2)\times8}{2}=48.
  \]
  \item 翻折后有宽\(8\)、高\(2\)的一段超出上边界，面积\(16\)。留在原长方形中的翻折部分面积\(48-16=32\)，且均与右侧未折纸重合，故双层面积为\(32\)。
  \item 也可看成直角边均为\(8\)的三角形，面积\(\dfrac{8\times8}{2}=32\)。
\end{enumerate}
\methodsummary{折叠越界割补。}
\end{solution}

\begin{solution}{15}{16}{8.0}{排列组合与概率}
\analysis{把圆圈在首尾相邻组处分类，然后分别转化成不相邻位置的链式选择。}
\begin{enumerate}[label=\textbf{Step \arabic*:},leftmargin=*]
  \item 若不取邻组\((8,1)\)，从依次排开的\(7\)个邻组中挑\(3\)个不相邻者，等价于从\(5\)个位置选\(3\)个，共\(\dbinom{5}{3}=10\)种。
  \item 若取了\((8,1)\)，另两组只能从\((2,3)\)至\((6,7)\)的\(5\)个邻组中不相邻地选出，等价于从\(4\)个位置选\(2\)个，共\(\dbinom{4}{2}=6\)种。
  \item 两类互斥且完备，总数\(10+6=16\)。编号不随圆环旋转重新识别。
\end{enumerate}
\methodsummary{圆环接缝与链式压缩。}
\end{solution}

\sectionbanner{三、综合题}{参考答案与解析}

\begin{solution}{16}{（1）23；（2）B}{4.0--6.0}{代数与方程应用}
\analysis{逐行扫描分割图，清点每行穿过的长方形种类数，再按\(X\times Y\)计分求和。}

\textbf{（1）}
第\(1\)行横穿标号\(3,4,6\)，共\(3\)种，得分\(1\times3=3\)；
第\(2\)行横穿\(5,2,6\)，共\(3\)种，得分\(2\times3=6\)；
第\(3\)行横穿\(7,6\)，共\(2\)种，得分\(3\times2=6\)；
第\(4\)行横穿\(8,1\)，共\(2\)种，得分\(4\times2=8\)。
总得分\(3+6+6+8=23\)。

\textbf{（2）}
第\(1\)行横穿\(8,7\)，共\(2\)种，得分\(1\times2=2\)；
第\(2\)行横穿\(8,5,4\)，共\(3\)种，得分\(2\times3=6\)；
第\(3\)行横穿\(8,2,1,3,4\)，共\(5\)种，得分\(3\times5=15\)；
第\(4\)行横穿\(8,2,6\)，共\(3\)种，得分\(4\times3=12\)。
总得分\(2+6+15+12=35\)，选B。

\methodsummary{逐行扫描、不重不漏地清点区块，再严格代入计分规则。}
\end{solution}

\begin{solution}{17}{（1）A；（2）B；（3）D}{4.0--8.0}{新定义与抽象建模}
\analysis{小\(A\)直接枚举；\(A=4\)用升降对称；插入最大元只做局部差量分析。}

\textbf{（1）}
\(A=3\)时全排列共\(6\)种。得分为\(1\)的有\(132,213,231,312\)，共\(4\)种，选A。

\textbf{（2）}
\(A=4\)时全排列\(24\)种。完全升序得\(3\)分、完全降序得\(0\)分各\(1\)种。将每个数\(k\)换为\(5-k\)后，升降关系全部翻转，故得\(m\)分与得\(3-m\)分的排列一一对应。于是得\(1\)分与得\(2\)分的种数相等，设为\(x\)，则
\[
1+x+x+1=24\implies x=11.
\]
故得\(2\)分的有\(11\)种，选B。

\textbf{（3）}
原数列\(1324\)基础得分为\(2\)。将\(5\)插入五个位置：
\begin{itemize}[leftmargin=*]
  \item \(51324\)：总分\(2\)；
  \item \(15324\)：总分\(2\)；
  \item \(13524\)：总分\(3\)；
  \item \(13254\)：总分\(2\)；
  \item \(13245\)：总分\(3\)。
\end{itemize}
可能得分只有\(2\)分和\(3\)分，选D。

\methodsummary{小数枚举、对称映射、插入时只看局部关系变化。}
\end{solution}

\begin{solution}{18}{（1）C；（2）5；（3）277}{6.0}{数论与整除逻辑}
\analysis{立方模\(9\)的余数只会循环出现\(\{0,1,8\}\)；三数和换负余数后落在\([-3,3]\)；四立方和从最大项降维试探。}

\textbf{（1）}
枚举可见\(a^3\bmod 9\)按\(1,8,0\)循环，只可能是\(\{0,1,8\}\)，不可能是\(4\)，选C。

\textbf{（2）}
每个余数属于\(\{0,1,8\}\)。把\(8\)看作\(-1\)，则三个余数之和\(S\in\{-3,-2,-1,0,1,2,3\}\)，对应模\(9\)余数为\(\{0,1,2,3,6,7,8\}\)，缺席\(4\)和\(5\)。故填\(5\)。

\textbf{（3）}
由\(d=11\)得
\[
a^3+b^3+c^3=2025-11^3=2025-1331=694,
\]
且\(a\le b\le c\le11\)。从大到小试\(c\)：\(c=9\)时\(729>694\)；\(c=8\)时两立方和\(182\)无法拼成；\(c=7\)时需\(351\)，而\(351-7^3=8=2^3\)，即\(a=2,b=7\)。又\(c\le6\)时三立方和至多\(648<694\)。验算
\[
2^3+7^3+7^3+11^3=8+343+343+1331=2025.
\]
故\(a,b,c\)为\(2,7,7\)，并列写作\(277\)。

\methodsummary{余数先找循环；多余数和用负余数缩范围；不定方程从最大项降维。}
\end{solution}
